\documentclass[11pt]{article}
\usepackage{mathblatt}
% Präambel des Setzers (werkzeuge/setzer-praeambel.tex, wortgleich); dazu \gteil (Muster T2U)
\makeatletter
\setlength{\mbpfbe}{0mm}% keine Punkte (pruefung.md § 5: nur beim Prüfstein)
\newlength{\gEin}\setlength{\gEin}{0pt}
\let\gSchrift\relax
\renewcommand{\pfkreuzzeile}[1]{\par\noindent\hangindent3.2em\hangafter1\hspace*{1.6em}$\square$\ #1\par}
\newcommand{\gkreuz}[1]{$\square$\ #1\hspace{2em}}
\newcommand{\gkreuzzeile}[2]{\par\noindent\hangindent2.2em\hangafter1\mbox{#1\,$\square$}\ #2\par}% Buchstabe, Option: eine je Zeile (neben Skizze)
% links, was man ansieht, rechts, was man tut; Spaltengrenze überall an derselben Stelle (gemeinsam.md „Satz“)
\newsavebox{\gbox}\newlength{\gLinks}
% v0.10 (Lehrer 13.10.): keine starre Halbierung – die linke Spalte ist so breit wie die Figur (+ 6 mm), die
% Antwort steht direkt daneben; ist die Figur breiter als 0,62 Zeile, steht die Antwort darunter
\newcommand{\gzwei}[2]{\par\smallskip\noindent\sbox{\gbox}{#1}%
  \setlength{\gLinks}{\dimexpr\wd\gbox+6mm\relax}%
  \ifdim\gLinks>0.62\linewidth
    \usebox{\gbox}\par\smallskip\noindent\begin{minipage}{\linewidth}#2\end{minipage}%
  \else
    \begin{minipage}[t]{\gLinks}\vspace{0pt}\usebox{\gbox}\end{minipage}%
    \begin{minipage}[t]{\dimexpr\linewidth-\gLinks\relax}\vspace{0pt}\raggedright #2\end{minipage}%
  \fi\par}
% v0.9 (12.10., pruefung.md § 5 „Marke mit B, Bedingungen rechts“): drittes Argument von pfaufg = Bedingungen
% der Prüfung (v0.10: \pfohnetr = TR durchgestrichen, „FOR“), klein grau im rechten Rand auf der Grundlinie der Nummer, Zeilen mit \\
\newcommand{\gbedrand}[2]{\ifblank{#2}{}{\rlap{\hspace*{\dimexpr#1+2mm\relax}%
  \parbox[t]{15mm}{\raggedright\scriptsize\leavevmode\color{mbgrau}\strut#2}}}}
% Teil eines Bündels (a, b …): Marke links im Rand, Bedingungen rechts im Rand, auf der Zeile des Teils
\newcommand{\gteilrand}[2]{\llap{\mbpfrandmarke{#1}\hspace*{\mbpfnr}}\gbedrand{\linewidth}{#2}}
% Teil eines Bündels: Buchstabe in eigener Spalte, Inhalt eingerückt (\gEin = Einzug, damit die Spaltengrenze von
% \gzwei an derselben Stelle bleibt)
\newenvironment{gbteil}[3]{\par\smallskip\noindent\gteilrand{#2}{#3}\makebox[\mbpfteil][l]{\textbf{#1}}%
  \begin{minipage}[t]{\dimexpr\linewidth-\mbpfteil\relax}\setlength{\gEin}{\mbpfteil}\raggedright\strut\ignorespaces}%
 {\end{minipage}\par}
% Grafik höchstens so breit wie die Spalte (Bündel nebeneinander, Koordinatensystem ohne Karo)
\newsavebox{\gpbox}
\newcommand{\gpassend}[1]{\sbox{\gpbox}{#1}\ifdim\wd\gpbox>\linewidth\typeout{GPASSEND: verkleinert \the\wd\gpbox\space auf \the\linewidth}\resizebox{\linewidth}{!}{\usebox{\gpbox}}\else\usebox{\gpbox}\fi}
% Skizzenfeld (v0.10, Lehrer 13.10.): Karo ab dem linken Rand über die Zeile, „Skizze“ klein grau oben links im Feld
% (wie \pfskizzenfeld der Vorlage); der rechte Rand bleibt den Bedingungen
\newcommand{\gskizze}[1]{\par\smallskip\noindent\begin{tikzpicture}
  \clip (0,0) rectangle ({\linewidth-1mm},{#1*10mm});
  \draw[black!16,line width=0.3pt,step=5mm] (0,0) grid ({\linewidth},{#1*10mm});
  \node[anchor=north west,inner sep=2pt,font=\scriptsize,text=mbgrau,fill=white] at (0,{#1*10mm}) {Skizze};
  \end{tikzpicture}\par}
\RenewDocumentEnvironment{pfaufg}{m m m}{%
  \begin{lrbox}{\mbaufgabebox}\begin{minipage}[t]{\linewidth}%
  \gSchrift\noindent\hspace*{\gEin}\mbpfrandmarke{#1}\gbedrand{\linewidth-\gEin}{#3}%
  \begin{minipage}[t]{\mbpfnr}\strut\textbf{#2}\end{minipage}%
  \begin{minipage}[t]{\dimexpr\linewidth-\gEin-\mbpfnr\relax}\raggedright\strut\ignorespaces}%
 {\end{minipage}%
  \end{minipage}\end{lrbox}%
  \setlength{\mbaufgabehoehe}{\dimexpr\ht\mbaufgabebox+\dp\mbaufgabebox\relax}%
  \ifdefempty{\gkopfnach}{\par\addvspace{5pt}\Needspace*{\mbaufgabehoehe}}{%
    % Stufenkopf, der mit dieser (ersten) Aufgabe zusammenbleibt: Platz für beide verlangen, dann den Kopf setzen
    % (12.10., Quadratische: ein Kopf stand allein unten, die Aufgabe mit Koordinatensystem auf der nächsten Seite)
    \par\Needspace*{\dimexpr\mbaufgabehoehe+40pt\relax}\gkopfnach\gdef\gkopfnach{}\par\addvspace{5pt}}%
  \noindent\usebox{\mbaufgabebox}\par\addvspace{5pt}}
% Stufenkopf vormerken: \gkopfvor{<Kopf>} – die nächste pfaufg setzt ihn, sobald beide zusammen Platz haben
\def\gkopfnach{}
\newcommand{\gkopfvor}[1]{\gdef\gkopfnach{#1}}
% Teilaufgabe mit Skizze links in gzwei: Buchstabe in eigener Spalte, darunter Text und Skizze
\newcommand{\gteil}[2]{\makebox[\mbpfteil][l]{\textbf{#1}}\parbox[t]{\dimexpr0.5\textwidth-\gEin-\mbpfnr-5mm-\mbpfteil\relax}{\raggedright #2}}
% Stufenkopf der Vorlage (mathblatt.sty 2026-10-12: Linie, Abstand, 13 pt, \pfjahre) ohne ihren Vorhalt von
% 14 Zeilen: der Setzer fordert vor dem Kopf 0,16 Seitenhöhe (Kopf + erste Aufgabe), mehr nicht
\setlength{\mbpfstufevorhalt}{0pt}
\makeatother
\newcommand{\karo}[1]{\par\smallskip\noindent\begin{tikzpicture}
  \clip (0,0) rectangle ({\linewidth-1mm},{#1*10mm});
  \draw[black!16,line width=0.3pt,step=5mm] (0,0) grid ({\linewidth},{#1*10mm});
  \end{tikzpicture}\par}
\newcommand{\kk}[1]{$\square$\,#1\hspace{0.9em}}
\newcommand{\linie}[1]{\,{\color{black!45}\rule[-1pt]{#1}{0.4pt}}\,}
\newcommand{\pkt}{\textperiodcentered{}}
% v0.14 (14.10.): Stufenkopf mit fünf Punkten (Option = Jahre für die Bestellkarte) oder „zuletzt 19“
\NewDocumentCommand{\gpunkte}{o m}{#2}
\newcommand{\gzuletzt}[1]{zuletzt #1}
% v0.7 (12.10.): Schülerfrage unter dem Stufenkopf; Ankreuzoption als Bild (Buchstabe und Kästchen darunter)
\newcommand{\gfrage}[1]{\par\nobreak\noindent\hspace*{\mbpfnr}{\small\itshape #1}\par\nobreak\addvspace{2pt}}
\newlength{\gbildbreite}\setlength{\gbildbreite}{0.32\linewidth}% Setzer setzt sie je Reihe (vier Bilder: 0,24)
\newcommand{\gbildoption}[3]{\begin{minipage}[t][#3][b]{\gbildbreite}\centering #2\par\smallskip\mbox{#1\,$\square$}\end{minipage}}
% Lineare (12.10.): Bild mit Buchstabe ohne Kästchen (Diagramm zuordnen); Aussage mit „wahr / falsch“ zum Ankreuzen
\newcommand{\gbildname}[3]{\begin{minipage}[t][#3][b]{\gbildbreite}\centering #2\par\smallskip\mbox{\textbf{#1}}\end{minipage}}
\newcommand{\gwahrfalsch}[2]{\par\noindent\hangindent2.2em\hangafter1\makebox[2.2em][l]{(#1)}#2\hfill\mbox{wahr $\square$\quad falsch $\square$}\par}
\makeatletter\setlength{\mbpfbe}{0mm}\makeatother
\begin{document}

\pfheftstil
\pflaufkopf[P10 \pkt{} mehr]{Brüche und Potenzen}{PZ10-1}\pfseitevon
\pfheftkopf{Brüche und Potenzen}{mehr}
\pfstufenzeile{Anteil gesucht \pkt{} Zahlen vergleichen \pkt{} Potenzen umwandeln}
% 2021-OS-B1b-m1 (Schwere I, rang 6)
\begin{pfaufg}{nach~P10\\’21~\pkt{}~1b}{1.}{}
Kennzeichne einen Anteil von $\tfrac{2}{5}$ in diesem Rechteck.
\par\smallskip\noindent\gpassend{\begin{tikzpicture}[x=0.45cm,y=0.45cm,line width=0.5pt]\draw[line width=0.4pt] (1,0) -- (1,2);\draw[line width=0.4pt] (2,0) -- (2,2);\draw[line width=0.4pt] (3,0) -- (3,2);\draw[line width=0.4pt] (4,0) -- (4,2);\draw[line width=0.4pt] (0,1) -- (5,1);\draw (0,0) rectangle (5,2);\end{tikzpicture}}\par
\karo{2}
\fusshilfe{\mbox{1:~4 Kästchen (beliebige Lage)}}
\end{pfaufg}
% 2019-OS-B1c-m1 (Schwere II, rang 7)
\begin{pfaufg}{nach~P10\\’19~\pkt{}~1c}{2.}{}
Kreuze an, welcher Anteil des Kreises grau ist:
\par\smallskip\noindent\gpassend{\begin{tikzpicture}[line width=0.6pt]\fill[black!50] (0,0) -- (0:0.8) arc (0:90:0.8) -- cycle;\fill[black!50] (0,0) -- (180:0.8) arc (180:225:0.8) -- cycle;\draw (0,0) circle (0.8);\draw (0,0) -- (0:0.8);\draw (0,0) -- (90:0.8);\draw (0,0) -- (180:0.8);\draw (0,0) -- (225:0.8);\draw (0,0) -- (270:0.8);\draw (0,0) -- (315:0.8);\end{tikzpicture}}\par
\par\smallskip\noindent\pfkreuz{A}{$\tfrac{2}{6}$}\pfkreuz{B}{$\tfrac{3}{8}$}\pfkreuz{C}{$\tfrac{2}{4}$}\pfkreuz{D}{$\tfrac{3}{6}$}
\fusshilfe{\mbox{2:~B}}
\end{pfaufg}
% 2017-OS-B1a-m1 (Schwere I, rang 8)
\begin{pfaufg}{nach~P10\\’17~\pkt{}~1a}{3.}{}
Schraffiere $\tfrac{4}{5}$ des Rechtecks.
\par\smallskip\noindent\gpassend{\begin{tikzpicture}[x=0.45cm,y=0.45cm,line width=0.5pt]\draw[line width=0.4pt] (1,0) -- (1,3);\draw[line width=0.4pt] (2,0) -- (2,3);\draw[line width=0.4pt] (3,0) -- (3,3);\draw[line width=0.4pt] (4,0) -- (4,3);\draw[line width=0.4pt] (0,1) -- (5,1);\draw[line width=0.4pt] (0,2) -- (5,2);\draw (0,0) rectangle (5,3);\end{tikzpicture}}\par
\karo{2}
\fusshilfe{\mbox{3:~12 Kästchen}}
\end{pfaufg}
% 2014-OS-B1i-m1 (Schwere I, rang 9)
\begin{pfaufg}{nach~P10\\’14~\pkt{}~1i}{4.}{}
Gib den Anteil der grauen Fläche als Bruch und in Prozent an.
\par\smallskip\noindent\gpassend{\begin{tikzpicture}[x=0.5cm,y=0.5cm,line width=0.5pt]\fill[black!25] (0,3) rectangle ++(1,1);\fill[black!25] (3,3) rectangle ++(1,1);\fill[black!25] (1,2) rectangle ++(1,1);\fill[black!25] (4,1) rectangle ++(1,1);\fill[black!25] (2,0) rectangle ++(1,1);\fill[black!25] (3,0) rectangle ++(1,1);\draw[line width=0.4pt] (1,0) -- (1,4);\draw[line width=0.4pt] (2,0) -- (2,4);\draw[line width=0.4pt] (3,0) -- (3,4);\draw[line width=0.4pt] (4,0) -- (4,4);\draw[line width=0.4pt] (0,1) -- (5,1);\draw[line width=0.4pt] (0,2) -- (5,2);\draw[line width=0.4pt] (0,3) -- (5,3);\draw (0,0) rectangle (5,4);\end{tikzpicture}}\par
\karo{4}
\fusshilfe{\mbox{4:~0,3; 30 \%}}
\end{pfaufg}
% 2018-OS-B1a-m1 (Schwere I, rang 10)
\begin{pfaufg}{nach~P10\\’18~\pkt{}~1a}{5.}{}
Bestimme $\tfrac{2}{3}$ von 1,5 kg.
\karo{3}
\fusshilfe{\mbox{5:~1 kg}}
\end{pfaufg}
% 2015-OS-B1h-m1 (Schwere II, rang 11)
\begin{pfaufg}{nach~P10\\’15~\pkt{}~1h}{6.}{}
Ein Regenfass fasst 900 Liter und ist zu $\tfrac{3}{4}$ gefüllt. Gib an, wie viel Liter noch hineinpassen, bis es voll ist.
\karo{2}
\fusshilfe{\mbox{6:~225 l}}
\end{pfaufg}
% 2019-OS-B1d-m1 (Schwere I, rang 14)
\begin{pfaufg}{nach~P10\\’19~\pkt{}~1d}{7.}{}
Unterstreiche die größte Zahl: $3^4$ · $9^3$ · $2^9$
\karo{4}
\fusshilfe{\mbox{7:~729}}
\end{pfaufg}
% 2017-OS-B1e-m1 (Schwere I, rang 15)
\begin{pfaufg}{nach~P10\\’17~\pkt{}~1e}{8.}{}
Kreuze die kleinste Zahl an:
\par\smallskip\noindent\pfkreuz{A}{$-$0,05}\pfkreuz{B}{$ - 5^3$}\pfkreuz{C}{$ - 5^2$}\pfkreuz{D}{$-$0,5}
\fusshilfe{\mbox{8:~B}}
\end{pfaufg}
% 2015-OS-B1c-m1 (Schwere II, rang 16)
\begin{pfaufg}{nach~P10\\’15~\pkt{}~1c}{9.}{}
Kreuze die wahre Aussage an:
\par\smallskip\noindent\pfkreuz{A}{$2{,}5 < \tfrac{5}{2}$}\pfkreuz{B}{$\tfrac{7}{4} > \tfrac{3}{2}$}\pfkreuz{C}{$\sqrt{3} > \tfrac{7}{4}$}
\fusshilfe{\mbox{9:~B}}
\end{pfaufg}
% 2014-OS-B1h-m1 (Schwere II, rang 17)
\begin{pfaufg}{nach~P10\\’14~\pkt{}~1h}{10.}{}
Ordne die Zahlen nach ihrer Größe, beginne mit der kleinsten: $ - \tfrac{1}{4}$ · 1,7 · $-$0,251 · $\sqrt{3}$
\karo{4}
\fusshilfe{\mbox{10:~$-$0,251 < $-\tfrac{1}{4}$ < 1,7 < √3}}
\end{pfaufg}
% 2020-OS-B1f-m1 (Schwere II, rang 18)
\begin{pfaufg}{nach~P10\\’20~\pkt{}~1f}{11.}{}
Gib die Zahl an, die genau in der Mitte zwischen $-$0,8 und $-$0,7 liegt.
\karo{2}
\fusshilfe{\mbox{11:~$-$0,75}}
\end{pfaufg}
% 2015-OS-B1b-m1 (Schwere I, rang 19)
\begin{pfaufg}{nach~P10\\’15~\pkt{}~1b}{12.}{}
Gib eine Zahl an, die kleiner als $-$80 ist.
\karo{2}
\fusshilfe{\mbox{12:~jede Zahl links von $-$80, z. B. $-$81}}
\end{pfaufg}
% 2014-OS-B1c-m1 (Schwere I, rang 20)
\begin{pfaufg}{nach~P10\\’14~\pkt{}~1c}{13.}{}
Gib eine Zahl an, die zwischen $\tfrac{1}{3}$ und $\tfrac{3}{5}$ liegt.
\karo{3}
\fusshilfe{\mbox{13:~z. B. 0,5}}
\end{pfaufg}
% 2020-OS-B1h-m1 (Schwere I, rang 23)
\begin{pfaufg}{nach~P10\\’20~\pkt{}~1h}{14.}{}
Gib die Zahl x an, für die die Gleichung $2^x = 32$ gilt.
\karo{3}
\fusshilfe{\mbox{14:~5}}
\end{pfaufg}
% 2017-OS-B1g-m1 (Schwere I, rang 24)
\begin{pfaufg}{nach~P10\\’17~\pkt{}~1g}{15.}{}
Schreibe die Zahl 10 000 000 als Zehnerpotenz.
\karo{3}
\fusshilfe{\mbox{15:~$10^{7}$}}
\end{pfaufg}
% 2016-OS-B1h-m1 (Schwere I, rang 25)
\begin{pfaufg}{nach~P10\\’16~\pkt{}~1h}{16.}{}
Gib $4{,}2 \cdot 10^6$ ohne abgetrennte Zehnerpotenz an.
\karo{3}
\fusshilfe{\mbox{16:~4 200 000}}
\end{pfaufg}
% 2019-OS-B1j-m1 (Schwere II, rang 26)
\begin{pfaufg}{nach~P10\\’19~\pkt{}~1j}{17.}{}
Gib als Dezimalzahl an: $3{,}4 \cdot 10^{ - 3}$
\karo{3}
\fusshilfe{\mbox{17:~0,0034}}
\end{pfaufg}
% 2015-OS-B1j-m1 (Schwere II, rang 27)
\begin{pfaufg}{nach~P10\\’15~\pkt{}~1j}{18.}{}
Kreuze die wahre Aussage an:
\par\smallskip\noindent\pfkreuz{A}{$\sqrt{( - 5)^2} =  - 5$}\pfkreuz{B}{$\sqrt{( - 5)^2} = +5$}\pfkreuz{C}{$\sqrt{( - 5)^2}$ ist nicht definiert.}
\fusshilfe{\mbox{18:~B}}
\end{pfaufg}
\par\vfill\noindent{\scriptsize\color{mbgrau}Originale: P10 2021 \pkt{} 1b, P10 2019 \pkt{} 1c, P10 2017 \pkt{} 1a, P10 2014 \pkt{} 1i, P10 2018 \pkt{} 1a, P10 2015 \pkt{} 1h, P10 2019 \pkt{} 1d, P10 2017 \pkt{} 1e, P10 2015 \pkt{} 1c, P10 2014 \pkt{} 1h, P10 2020 \pkt{} 1f, P10 2015 \pkt{} 1b, P10 2014 \pkt{} 1c, P10 2020 \pkt{} 1h, P10 2017 \pkt{} 1g, P10 2016 \pkt{} 1h, P10 2019 \pkt{} 1j, P10 2015 \pkt{} 1j\par}
\end{document}